\section{Case studies}
\label{sec:case-studies}

This section reports two research languages for scientific computing built on \PurePy: \Fortl and \Fluid. \todo{Is there a way to do this without deanonymising the paper?}

\subsection{\Fortl}
\label{sec:case-study-fortl}

\Fortl is a statically-typed research language for scientific and numerical computing, with a Python-like
syntax. Its distinctive feature is a system of graded numerical types: types are indexed and carry structure,
so domain concepts such as units of measure and physical quantities are expressed in the types and enforced by
the type checker. A length divided by a time is typed as a speed, and a program that mixes incompatible units
is rejected. \Fortl aims to let code state the meaning, intent, and limitations of a scientific model, and to
make its predictions reproducible and connected to their abstract description. \todo{Confirm and describe the
relationship to \PurePy}

\subsection{\Fluid}
\label{sec:case-study-fluid}

\Fluid is a pure, Pythonic, provenance-aware functional language for transparent research outputs. It is built
on \PurePy, adding a runtime that records fine-grained dependencies between the inputs and outputs of a
computation, so that a figure or summary can be linked back to the underlying data. \Fluid implements the
partial-application and piecewise extensions of \secref{extensions}.
